\documentclass[11pt,a4paper]{article}

\usepackage[T1]{fontenc}
\usepackage[utf8]{inputenc}
\usepackage{lmodern}
\usepackage[margin=2.6cm]{geometry}
\usepackage{microtype}
\usepackage{amsmath,amssymb,amsthm}
\usepackage{booktabs}
\usepackage{longtable}
\usepackage{tabularx}
\usepackage{array}
\usepackage{enumitem}
\usepackage{xcolor}
\usepackage{fancyvrb}
\usepackage{algorithm}
\usepackage{algpseudocode}
\usepackage[hidelinks]{hyperref}

\newtheorem{proposition}{Proposition}
\theoremstyle{definition}
\newtheorem{definition}{Definition}
\theoremstyle{remark}
\newtheorem{remark}{Remark}

\newcommand{\code}[1]{\texttt{#1}}
% A claim to be confirmed: the text in red, the reason in a footnote.
\newcommand{\tocheck}[2]{\textcolor{red}{#1}\footnote{\textcolor{red}{To check: #2}}}
\newcommand{\cl}{\mathcal{L}}
\newcommand{\pos}{^{+}}
\newcommand{\nega}{^{-}}
\newcommand{\op}{^{\mathrm{op}}}
\newcommand{\Eq}[1]{\mathrm{Eq}#1}
\newcommand{\Th}{\operatorname{Th}}

\fvset{fontsize=\small,frame=single,framesep=2mm}
\setlist{itemsep=2pt,topsep=4pt}

\title{Conceptual Exploration of Equational Theories of Magmas\\[4pt]
\large Implementation notes, two case studies, and a user's guide}
\author{The \code{conceptual-exploration} library}
\date{October 2026}

\begin{document}
\maketitle

\begin{abstract}
This document describes how the \code{conceptual-exploration} library carries out the attribute exploration of implications between equational laws of magmas, the question studied by the Equational Theories Project. It is written for a reader who knows attribute exploration as presented in Ganter and Obiedkov's \emph{Conceptual Exploration}, and it concentrates on what is done differently from the book: partial objects, symmetries, experts that cannot always decide, and the resumption of an exploration when an implication accepted without proof is refuted. It then describes what is specific to equational theories: the attributes and objects, background knowledge derived by bounded equational reasoning, and three experts: one searching for finite counterexamples with the Z3 SMT solver, one using the Vampire theorem prover, which also refutes implications with free magmas, and one deciding questions by Knuth--Bendix completion. Three case studies are discussed in detail: sixteen classical laws, all thirty-one one-variable laws with at most four operations, and the ETP laws of order at most two and three, whose implications from several laws have been determined, or almost, by other means in other work. A user's guide closes the document.
\end{abstract}

\tableofcontents

%======================================================================
\section{Introduction}
%======================================================================

A \emph{magma} is a set $G$ with a binary operation $*$. An \emph{equational law} (or \emph{law}) is an equation $s = t$ between two terms built from variables with $*$, read as universally quantified: $G$ \emph{satisfies} $s = t$, written $G \models s = t$, if the two sides agree under every assignment of elements of $G$ to the variables. Examples are commutativity $x * y = y * x$, associativity $(x * y) * z = x * (y * z)$ and idempotence $x = x * x$. The \emph{order} of a law is the number of occurrences of $*$ in its two sides together.

The Equational Theories Project (ETP)~\cite{ETP25}, initiated by Terence Tao in 2024, set out to decide, for every pair of laws $E, F$ from a list of 4694 laws of order at most four, whether every magma satisfying $E$ also satisfies $F$. Its laws are numbered, and we keep its numbers: $\Eq{43}$ is commutativity, $\Eq{4512}$ associativity, and so on.

From the point of view of Formal Concept Analysis, the natural object of study is the formal context
\[
  \mathbb{K}_{\mathrm{Mag}} = (\mathcal{G}, M, \models),
\]
where $M$ is a finite set of laws and $\mathcal{G}$ is the class of all magmas. Its object class is not even a set, but its implications are perfectly well defined: $P \to C$ holds in $\mathbb{K}_{\mathrm{Mag}}$ if every magma satisfying all laws of $P$ satisfies all laws of $C$. By Birkhoff's completeness theorem, this is the case exactly when every law of $C$ can be derived from the laws of $P$ by the rules of equational logic. Its intents are the sets $\Th(G) \cap M$ for magmas $G$, the traces on $M$ of equational theories. Attribute exploration is the natural tool for computing the canonical basis of $\mathbb{K}_{\mathrm{Mag}}$ and, with it, all its intents: an expert either confirms an implication, ideally by a proof, or refutes it by a magma.

The library implements the exploration algorithm of the book~\cite{GO16} in a general form, in the Python package \code{conceptual\_exploration}, and the equational-theories domain in the directory \path{explorations/equational_theories}. Section~\ref{sec:library} describes how the general algorithm differs from the book. Section~\ref{sec:et} describes the domain: attributes, objects, the duality symmetry, background knowledge, and the experts. Sections~\ref{sec:case16}, \ref{sec:case31}, \ref{sec:case3} and~\ref{sec:semigroups} discuss four explorations, and Section~\ref{sec:guide} explains how to run one's own.

%======================================================================
\section{Attribute exploration as implemented}\label{sec:library}
%======================================================================

\subsection{Overview}

The exploration loop is the one of the book: the premises are the sets closed under the implications known so far, enumerated with NextClosure in lectic order, and for each premise $P$ the expert is asked about $P \to C$, where $C$ is what the current context suggests $P$ implies. Four features change what happens around this loop:
\begin{enumerate}
  \item objects are \emph{partial}: for some attributes it may be unknown whether an object has them (Section~\ref{sec:partial});
  \item a group of \emph{symmetries} acts on the attributes, and implications and objects are added together with their images (Section~\ref{sec:symmetry});
  \item the expert may \emph{fail to decide} a question, in which case parts of the conclusion may be accepted without proof or rejected (Section~\ref{sec:undecided});
  \item when an implication accepted without proof is later refuted, the exploration is \emph{resumed} at the point where it went wrong (Section~\ref{sec:rounds}).
\end{enumerate}
Algorithm~\ref{alg:exploration} summarizes the loop; the notation is explained in the following subsections.

\begin{algorithm}[t]
\caption{The exploration loop (\code{AttributeExploration.run})}\label{alg:exploration}
\begin{algorithmic}[1]
\Require attributes $M$ with a linear order, background implications $\mathcal{B}$, symmetries $\Phi$, expert
\State $\cl \gets \mathcal{B}$ closed under $\Phi$; \quad $\mathcal{T} \gets \cl$ \Comment{all implications; trusted ones}
\State $P \gets \cl(\emptyset)$
\Loop
  \State $C \gets P^{\diamond} \setminus P$ \Comment{context closure, Section~\ref{sec:partial}}
  \If{$C = \emptyset$} \State $P \gets$ next $\cl$-closed set after $P$; \textbf{continue} \EndIf
  \State ask the expert about $P \to C$
  \If{the answer is an object $g$ refuting $P \to C$}
    \State add $g$ and its images under $\Phi$, completed with $\mathcal{T}$; \textbf{continue} with the same $P$
  \ElsIf{the answer is ``accept''}
    \State add $P \to C$ and its images to $\cl$ (and to $\mathcal{T}$ if the expert decided it)
  \Else \Comment{undecided: proved part $D \subseteq C$, partial object, flag}
    \State add the partial object; add $P \to D$ to $\cl$ and $\mathcal{T}$
    \State add $P \to C \setminus D$ to $\cl$ only (unconfirmed), or reject it
  \EndIf
  \If{an object violates an implication of $\cl \setminus \mathcal{T}$} \State \textbf{stop} and resume (Section~\ref{sec:rounds}) \EndIf
  \State $P \gets$ next $\cl$-closed set after $P$ (or stop if there is none)
\EndLoop
\end{algorithmic}
\end{algorithm}

\subsection{Partial objects and the question asked}\label{sec:partial}

An object $g$ of the context is a \emph{partial object}: a pair $(g\pos, g\nega)$ of disjoint sets of attributes, those $g$ is known to have and those it is known not to have. The attributes in neither set are unknown. An object refutes $P \to C$ if $P \subseteq g\pos$ and $C \cap g\nega \neq \emptyset$.

The conclusion offered to the expert is computed with the closure
\[
  P^{\diamond} = M \setminus \bigcup \{\, g\nega : g \text{ an object with } P \subseteq g\pos \,\}.
\]
An attribute belongs to $P^{\diamond}$ unless some object known to have all of $P$ is known not to have it; $P^{\diamond}$ is thus the largest conclusion not refuted by the data at hand. When every row is complete, $P^{\diamond} = P''$ and the question is the one of the book. It is the cautious choice for incomplete data, and it is what makes partial rows useful: an object whose row is only partly known still removes from later conclusions the attributes it is known to lack.

In the equational-theories exploration, finite magmas always have complete rows (the expert evaluates every law on every counterexample), while free magmas (Section~\ref{sec:free}) may have unknown entries.

\paragraph{Completion.}
Every object is \emph{completed} with the \emph{trusted} implications $\mathcal{T}$ (background, confirmed, and their images; see Section~\ref{sec:undecided}): its positive side is replaced by its $\mathcal{T}$-closure, and every unknown attribute $m$ whose addition would force an attribute of $g\nega$ is added to the negative side:
\[
  g\pos := \mathcal{T}(g\pos), \qquad
  g\nega := g\nega \cup \{\, m : \mathcal{T}(g\pos \cup \{m\}) \cap g\nega \neq \emptyset \,\},
\]
the second step repeated until nothing changes. If the closure of the positive side meets the negative side, the object contradicts a trusted implication, and the program stops with an error: such an object can only come from a faulty expert. Objects are completed when they are added, and all objects are brought up to date whenever new trusted implications are accepted. Completion is what eventually fills in some unknown entries of partial rows.

\subsection{Background knowledge}

Background implications are given before the exploration starts, as in the book. They take part in the closure of premises and in the completion of objects, and they are reported with the source \code{background}. Without the features of the next subsections, the accepted implications form the canonical basis of the context relative to the background implications.

\subsection{Symmetries}\label{sec:symmetry}

A \emph{symmetry} (``mapping'' in the code) is a map $\varphi$ from the attributes to the attributes. When an implication $P \to C$ is accepted, so is its image $\varphi(P) \to \varphi(C)$ for every $\varphi \in \Phi$, unless the theory already entails it; images are reported with the source \code{mapped}, or \code{mapped unconfirmed} for images of unconfirmed implications. When an object $g$ is added, so is, for every $\varphi \in \Phi$, a copy of $g$ with the row
\[
  \bigl(\varphi^{-1}(g\pos), \; \varphi^{-1}(g\nega)\bigr),
\]
the object that has an attribute $m$ exactly when $g$ has $\varphi(m)$. Copies are named $g\#k$, where $k$ is the number of the mapping (the original is $g\#0$), and a copy whose row repeats one already added is skipped. Each implication and each object is mapped by each $\varphi$ once and never by a composition, so $\Phi$ should be closed under composition, for instance a group. (In general a map may also send an attribute to a truth value; this is used in rule exploration, not here.)

Since the theory is closed under $\Phi$ at all times, the premise $\varphi(P)$ of the image of an accepted implication is not closed and never becomes a question: the expert is asked about one implication per orbit. The images of an unconfirmed implication are no more certain than it, but they steer the enumeration as it does. With \code{map\_unconfirmed=False} (\code{-{}-ask-unconfirmed-duals} in \code{explore.py}), they are added only once the implication is confirmed; until then the expert is asked about them when the enumeration reaches their premises, and if it proves an image, the implication is recorded as confirmed as well. The price is that the resulting set of implications, background, accepted and mapped together, is only a \emph{cover} of the theory: a set of implications that generates it, possibly with redundant implications, and in general not its canonical basis.

\subsection{Experts that cannot always decide}\label{sec:undecided}

The book assumes an expert who always answers correctly. The experts used here are programs with time limits, and a question may be left undecided. The library distinguishes five sources of implications:
\begin{description}
  \item[\code{background}] given in advance;
  \item[\code{confirmed}] accepted by an expert who decided the question;
  \item[\code{unconfirmed}] accepted without a decision: the expert found no counterexample but cannot rule one out;
  \item[\code{mapped}] the image of a background or confirmed implication under a symmetry;
  \item[\code{mapped unconfirmed}] the image of an unconfirmed implication, no more certain than it; it becomes \code{mapped} if the implication is confirmed later.
\end{description}
An expert reports that it has not decided a question it accepts through \code{is\_conclusive()}. In addition, an expert may answer with an \code{Undecided} value, which carries
\begin{itemize}
  \item a set $D \subseteq C$ of conclusion attributes it did prove: $P \to D$ is accepted as confirmed;
  \item optionally a partial object, which must not refute $P \to C$ and is added to the context (its unknown entries may be filled later by completion);
  \item a flag telling what to do with the rest $P \to C \setminus D$: it is accepted as \emph{unconfirmed} if the expert knows nothing against it, or \emph{rejected} if the expert knows it to be false without knowing which attribute fails.
\end{itemize}

\paragraph{Trusted and untrusted implications.}
Unconfirmed implications take part in the closure that generates the premises: they steer the enumeration as accepted implications do. They are, however, not used to complete objects, since they may be false and would then write false values into rows. Every object added later is checked against them; an object violating an unconfirmed implication refutes it.

\paragraph{Rejected implications.}
A rejected implication is not accepted, and the enumeration continues with the next closed set as if the context had no question about $P$. The theory remains sound. It may, however, be incomplete at $P$: the expert knows that $P \to C \setminus D$ fails as a whole, but some of the attributes of $C \setminus D$ may still follow from $P$. Rejected implications are listed separately in the results. If the rest has a single attribute $c$, though, it is not rejected: some object has the attributes of $P$, and those of $D$, which $P$ implies, and lacks $c$, and nothing else is known about it. Such a \emph{witness} is a counterexample, if a partial one, and it joins the context with the row $(P \cup D, \{c\})$; in the results, witnesses are named \code{W}$k$.

A rejected implication is a fact, which no row can express: it says that some attribute of $C \setminus D$ fails, but not which. It is used as a constraint instead. An unconfirmed implication that, together with the theory and with its images under the symmetries if they are added, would entail a rejected implication is not accepted; it is reported as \emph{conflicting}, together with the rejected implication. A rejected implication is never entailed when it is recorded, since its premise is closed under the theory; it can become entailed only by a later unconfirmed implication, through its images or in a later round (Section~\ref{sec:rounds}), so that this check suffices.

\subsection{Resuming after a refuted unconfirmed implication}\label{sec:rounds}

When an object refutes an unconfirmed implication $P \to C$, all closed sets enumerated after $P$ were computed with a false implication, and the enumeration from $P$ on cannot be trusted. The exploration then stops and is resumed (\code{explore\_in\_rounds}): a new round starts with
\begin{itemize}
  \item the implications accepted before $P \to C$, whatever their source;
  \item the \emph{confirmed} implications accepted after it, which are true whatever steered the enumeration to them;
  \item every object found so far, and every rejected implication, which is false whatever steered the enumeration to it;
\end{itemize}
and resumes NextClosure at $P$. Since NextClosure can continue from any set, and every closed set lectically before $P$ was enumerated without the false implication, nothing before $P$ is repeated. The false implication and the unconfirmed implications accepted after it are dropped; their questions are asked again if they arise. The implications found are then a cover of the theory, sound and complete up to rejected implications, but, because the confirmed implications kept were found under a different theory, not necessarily its canonical basis.

\subsection{Further details}\label{sec:details}

\begin{itemize}
  \item \textbf{Several questions per premise.} After a counterexample, the context closure of the same premise is recomputed and, if it is still larger than the premise, the expert is asked again, as in the book.
  \item \textbf{Continuing after an acceptance.} After $P \to C$ is accepted, the next closed set is computed from $P$ under the extended theory. NextClosure reuses the step that produced $P$: the closure of $P$ under the new theory is the next closed set if it adds no attribute lectically smaller than the one just added.
  \item \textbf{Rows of counterexamples.} With \code{evaluate\_all}, the expert is given all attributes and returns the full row of a counterexample, not only its values on the question's attributes. The equational-theories exploration always does this.
  \item \textbf{Bit sets.} Attribute sets, rows and implications are stored as bit masks, and distinct rows are stored once with the objects sharing them; the context closure and the implication closure (a variant of Wild's algorithm) are integer operations.
  \item \textbf{Reading the result.} For display, \code{simplified\_rules} simplifies the accepted implications into \emph{rules}: each premise is cut to the attributes the rest of it does not entail, and each conclusion to the attributes that do not follow from the premise by way of the background, the other rules and all their images. A rule left with no conclusion follows from the others and is dropped, so that the rules are also reduced: none of them follows from the others together with the background and the symmetries. The rules thus form a basis of the theory relative to the background and the symmetries, though not in general its canonical basis, nor one with as few implications as possible. A premise that entails every attribute is shown as concluding \emph{everything}. The rules, read together with the background and the symmetries, entail exactly what the accepted implications do.
\end{itemize}

%======================================================================
\section{Exploring equational theories}\label{sec:et}
%======================================================================

\subsection{Attributes}

The attributes are the laws listed in a JSON file, each with an equation string and optionally a name and an ETP number:
\begin{Verbatim}
[
  {"id": 43, "name": "Commutativity", "equation": "(x * y) = (y * x)"},
  {"id": 4512, "name": "Associativity", "equation": "((x * y) * z) = (x * (y * z))"},
  {"name": "Medial / Entropic", "equation": "((x * y) * (z * w)) = ((x * z) * (y * w))"}
]
\end{Verbatim}
Terms are parsed into trees with variables as leaves and $*$ as the only operation; laws are compared syntactically. The set of laws must be closed under duality (Section~\ref{sec:duality}).

\subsection{Objects: finite and free magmas}

A \emph{finite magma} is given by its Cayley table on $\{0, \dots, n-1\}$. Whether it satisfies a law is decided by evaluating both sides under all $n^k$ assignments to the $k$ variables of the law, so rows of finite magmas are always complete.

A \emph{free magma} (Section~\ref{sec:free}) is a genuine magma, but usually an infinite one, so the library does not construct it; it represents it by its premise and computes only its row, which may be partial.

\subsection{Duality}\label{sec:duality}

The \emph{dual} of a term reverses all products, $(s * t)\op = t\op * s\op$, and the dual of a law $s = t$ is $s\op = t\op$. The dual $G\op$ of a magma has the transposed table, $a *\op b = b * a$, and $G\op \models E\op$ if and only if $G \models E$. The map $E \mapsto E\op$, with laws identified up to renaming variables and swapping sides, is the only symmetry used; it is an involution, so $\Phi = \{\,{}\op\}$ is closed under composition. Counterexamples are thus added together with their duals, named \code{M}$k$ and \code{M}$k$\code{\^{}op} in the results, and implications together with their dual implications.

When every object satisfies associativity (option \code{-{}-axiom}, which adds the empty-premise implication concluding the given laws to the background and drops seeded magmas that violate them), the objects are semigroups, and a law is determined by the two words obtained by dropping the parentheses. Laws are then compared as words, so that the dual of a law is the law with both words reversed, again up to renaming variables and swapping sides; the dual of a semigroup is a semigroup. Option \code{-{}-no-duality} explores without the symmetry, for sets of laws that are not closed under it.

\subsection{Background implications from bounded equational reasoning}\label{sec:background}

For laws in a single variable, the library can derive background implications syntactically. A derivation uses three rules:
\begin{enumerate}
  \item multiply both sides by the same term from the same side: from $s = t$ derive $u * s = u * t$ and $s * u = t * u$;
  \item substitute a term for the variable: from $s = t$ derive $s[x := u] = t[x := u]$;
  \item replace one occurrence of a subterm by an equal one: from $a = b$ and an equation $e$ derive $e$ with one occurrence of $a$ replaced by $b$ (or vice versa). This includes symmetry and transitivity.
\end{enumerate}
\tocheck{These rules are complete for one-variable equational logic}{every equational consequence of one-variable laws should be derivable as a chain of replacements of instances $u[x := v] \to w[x := v]$ of premise laws inside contexts, which rules~2, 1 and~3 provide in this order; this argument has not been written out in full.}, but derivations may pass through arbitrarily large equations. Derivations are therefore bounded by an \emph{order} $N$: all equations they pass through have order at most $N$. Within this bound, all rule instances are computed once, as Horn clauses between numbered equations, and the laws a set of laws derives are obtained by forward chaining.

A \emph{transient order} $N' > N$ admits derivations of one more shape: a derived equation is multiplied or substituted into once, giving an equation of order at most $N'$, and both sides of it are then rewritten independently with the derived equations, each rewrite replacing a side of a derived equation by the other side when that is no larger. Every pair of terms reached from the two sides that forms an equation of order at most $N$ is derived. For example, substituting $x * x$ for $x$ in $x * x = ((x * x) * x) * x$ gives an equation of order~10, which $x = (x * x) * (x * x)$ rewrites to $x = (x * (x * x)) * (x * x)$, of order~4. Such steps alternate with the forward chaining until nothing new is derived.

The derivation closure $\sigma$ is a closure operator on sets of laws, and the background implications are computed from it in one of two ways:
\begin{itemize}
  \item with a \emph{premise size} $K$ (default $2$), every set $P$ of at most $K$ laws yields $P \to \sigma(P) \setminus P$, unless the implications kept so far entail it. The result is sound but may miss implications that need larger premises;
  \item with premise size $0$, the canonical basis of $\sigma$ is computed with NextClosure, which visits every $\sigma$-closed set and is much slower.
\end{itemize}
In both cases the background is sound: every background implication is valid in all magmas. Laws with several variables are not supported, since rule~2 would then have to rename variables apart and identify laws up to renaming.

\subsection{Background implications from the ETP}\label{sec:etp}

The ETP has decided every implication between two of its 4694 laws, and publishes the outcomes as an implication graph\footnote{\url{https://teorth.github.io/equational_theories/implications/graph.json}}: a $4694 \times 4694$ array, row the premise and column the conclusion, run-length encoded. Together with the ETP's list of laws\footnote{\url{https://github.com/teorth/equational_theories/blob/main/data/equations.txt}}, it gives background implications for laws with any number of variables:
\begin{itemize}
  \item every explored law is matched with an ETP law by a canonical form: its variables are renamed in the order of their first occurrence, for both orders of the two sides, and the smaller of the two strings is taken. Two laws have the same canonical form exactly when one becomes the other by renaming variables and swapping sides. A law whose \code{id} is not the ETP's number of the matched law stops the program with an error; laws the ETP does not list take no part;
  \item the laws that $x = x$ implies in the graph, which hold in every magma, form the conclusion of an implication with empty premise;
  \item every other listed law $E$ yields $E \to \{F : E \text{ implies } F\}$, leaving out the laws valid in every magma.
\end{itemize}
The exploration then only asks about premises of several laws, since every implication with a single premise between listed laws is entailed by the background. \tocheck{The ETP's implications are proved in its Lean formalization}{the graph's outcomes 3 and 7 (implication true) and 2 and 6 (false) are read as ``explicitly or implicitly proved''; whether every such outcome is backed by a Lean proof was not checked.}, so this background is as trustworthy as the ETP; it is reported with the source \code{background}, like derived implications, and both kinds can be used together.

The graph also gives the implications that the ETP refutes, though not its counterexamples, some of which are infinite. With \code{-{}-etp-refutations}, each refuted implication $E \to F$ between listed laws is added to the context as a partial object (Section~\ref{sec:partial}): the ETP's counterexample satisfies $E$ and every law that $E$ implies, and violates $F$ and every law that implies $F$, and nothing else is known about it. These objects are reduced together with the seeded magmas (Section~\ref{sec:seeding}), so that only refutations that no other object accounts for remain; they are named \code{R}$k$ in the results. They matter when the ETP's counterexample is infinite, or finite but beyond the reach of the expert, since a question refuted only by such a magma would otherwise be accepted as unconfirmed.

The ETP data used in this document are those of the ETP repository at commit \code{1aec8a7}, whose \code{equations.txt} and \code{graph.json} are recorded with their checksums in \code{etp\_version.json}. The script \code{fetch\_etp.py} downloads them, and the ETP's files of finite magmas used in Section~\ref{sec:case3}, which are recorded there as well, and checks each file against its checksum: \code{fetch\_etp.py DIR} puts \code{equations.txt} and \code{graph.json} into \code{DIR}, and the files of magmas into \code{DIR/magmas}.

The script \code{etp.py} writes chosen ETP laws, together with their duals, as an equations file with the ETP's numbers as ids, so that the laws to explore can be chosen by number.

\subsection{The Z3 expert}\label{sec:z3}

The \code{MagmaExpert} answers a question $P \to C$ as follows.
\begin{enumerate}
  \item It checks a \emph{pool} of magmas: 24 small standard magmas (projections, constants, cyclic groups, the two-element Boolean magmas, a rock--paper--scissors tournament, all two-element tables), every counterexample found so far, and any magmas given as seeds. The first magma satisfying $P$ and violating a law of $C$ is the counterexample.
  \item Otherwise, for each size $n$ from a minimum to a maximum (by default $1$ to $6$), it asks the Z3 SMT solver for a table of size $n$ satisfying $P$ and violating some law of $C$. The operation is an uninterpreted function restricted to $\{0, \dots, n-1\}$, and the laws are grounded over all assignments. The first table found is the counterexample and joins the pool.
  \item If no size yields a table, the implication is accepted. It is reported as confirmed if Z3 answered ``unsatisfiable'' for every size and as unconfirmed if it ran out of time (by default 30 seconds) for some size.
\end{enumerate}
A ``confirmed'' implication thus only means that it has no counterexample of the searched sizes. It may be false, refuted only by a larger or an infinite magma, and both case studies show that this happens.

\subsection{The Vampire expert}\label{sec:vampire}

The \code{VampireExpert} uses the first-order theorem prover Vampire. A question $P \to C$ becomes a TPTP problem whose axioms are the universal closures of the laws of $P$ and whose conjecture is the conjunction of the laws of $C$. Vampire then answers with a status: \emph{Theorem} (it found a proof), \emph{CounterSatisfiable} (the axioms do not imply the conjecture, shown by a finite model or by a saturated clause set), or a timeout.

The expert proceeds through the following steps and stops at the first one that decides the question.
\begin{enumerate}
  \item \textbf{Pool and memory.} The pool is checked as by the Z3 expert. A question entailed by an implication proved before ($P' \subseteq P$, $C \subseteq C'$) is accepted, and an undecided answer is reused if the same question is asked again.
  \item \textbf{Proof.} Vampire's portfolio mode (\code{-{}-mode casc}) runs on the problem, on all cores, since nothing else runs at this point (option \code{-{}-vampire-first-time-limit}, by default the proving time limit of 10 seconds). A proof confirms the implication. If the run instead ends with a finite model, the model is the counterexample. On the 100 whole conclusions of the benchmark described in Section~\ref{sec:free}, eight cores decided the same 81 as one core and none in more than 1.8 seconds, against 9.6 seconds for the slowest decision on one core; with a limit of 3 seconds, the step took 70\% less time than on one core with 10 seconds, and lost no decision.
  \item \textbf{Quick model search.} Vampire's finite model builder (\code{-{}-saturation\_algorithm fmb}) looks for a counterexample of any size (by default for 1 second).
  \item \textbf{One law shown false.} If $C$ has a single law and step~2 ended with CounterSatisfiable, the free magma over $P$ violates that law and refutes the implication. No further call is made: the rest of its row is decided by the context already, the laws outside $P \cup C$ being denied for $P$ by its objects.
  \item \textbf{Free magma.} The row of the free magma over $P$ is computed, as described in Section~\ref{sec:free}. If a law of $C$ is on its negative side, the free magma refutes the implication. If every law of $C$ is on its positive side, the implication is proved. If step~2 ended with CounterSatisfiable and a single law $c$ of $C$ is not on the positive side, the free magma violates $c$: a brief search (1~second) looks for a finite model violating $c$, and the free magma refutes the implication if it finds none.
  \item \textbf{Per-law model search.} For each law $c$ of $C$ that the free magma did not prove, the model builder looks for a model of $P$ violating $c$ (by default for 10 seconds each), all searches sharing the cores. This is much faster than a search for a model violating the conjunction $C$, which has to consider all the ways of violating it at once. Optionally, the model finder Mace4\footnote{\url{https://www.cs.unm.edu/~mccune/prover9/}} searches as well (option \code{-{}-mace4-time-limit}, off by default), with its own propagation and symmetry reduction, which sometimes finds models that Vampire's model builder does not reach in time. Each Mace4 search is for one law at one size, and the sizes are swept in parallel: a core that finishes a search, having found a model or shown that there is none of that size, takes the next search, for the smallest size not yet started; the sweep gets the given time in all. A search at a fixed size where a model exists is often fast, while ruling out the smaller sizes first, as a search size by size has to, can take much longer. Each model found is checked and joins the pool, and the searches for the laws it violates stop; the others go on until every law of $C$ is violated or the time is up. The smallest model found refutes $P \to C$, and the pool then answers the next questions about $P$, whose conclusions are the laws the model satisfies, without a further call. A law whose search ran out of time is not searched for again with the same premise. With option \code{-{}-model-premise simplified}, the searches use a subset of $P$ whose closure under the background and the implications proved so far is $P$, keeping the laws that every object satisfies, which the background assumes rather than derives; \code{both} searches with both premises. Every model is checked against the whole premise.
  \item \textbf{Undecided.} Otherwise, the laws of $C$ that the free magma proved are accepted as confirmed. The rest of $C$, which has at least two laws if step~2 ended with CounterSatisfiable, is rejected if it did, that is, if Vampire showed $P \to C$ false as a whole, and is accepted as unconfirmed otherwise. The partial row of the free magma is not added to the context. Its positive side holds in every object of the context that satisfies $P$, since objects are completed with the trusted implications, and every law on its negative side lies outside $C$ and is therefore denied by some object satisfying $P$; so the row could refute nothing that the context does not already refute.
\end{enumerate}
Every finite model is read from the output of Vampire or Mace4 and checked against the laws; a model that does not refute the question stops the program with an error. Vampire's proofs and its CounterSatisfiable verdicts obtained by saturation are trusted. The expert counts its answers by kind (\code{pool}, \code{proved}, \code{model}, \code{free magma}, \code{rejected}, \code{unconfirmed}), and the counts are printed at the end.

Unlike the Z3 expert, the Vampire expert never accepts an implication as confirmed without a proof, and it never accepts an implication it knows to be false.

\subsection{Free magmas}\label{sec:free}

Let $P$ be a set of laws. Call two terms $s$ and $t$, in countably many variables, \emph{equivalent} if $P$ derives $s = t$. The \emph{free magma} $F_P$ over $P$ is the set of equivalence classes, with the operation $[s] * [t] = [s * t]$, which is well defined because $P$ derives $s * t = s' * t'$ whenever it derives $s = s'$ and $t = t'$. It satisfies $P$, and by Birkhoff's theorem it satisfies a law $L$ exactly when $P$ derives $L$. Hence:
\begin{proposition}
For every finite set $M$ of laws and every $P \subseteq M$, the row of $F_P$ is $\{\, L \in M : P \vdash L \,\}$, the closure of $P$ in $\mathbb{K}_{\mathrm{Mag}}$. In particular, $F_P$ refutes every false implication with premise $P$.
\end{proposition}
For one-variable laws, the submagma of $F_P$ generated by one variable already has the same row. $F_P$ is usually infinite, so it has no table. Its row is computed law by law. The laws that $P$ implies by the background implications and by the implications Vampire has proved so far are on the positive side without asking. For each other law $L$, Vampire runs in portfolio mode on the problem ``$P$ implies $L$'', all laws in parallel and the laws of the conclusion first. A proof puts $L$ on the positive side, CounterSatisfiable on the negative side, and a timeout leaves $L$ unknown. With \code{-{}-eprover-time-limit}, the E prover\footnote{\url{https://eprover.org}} decides each law of the conclusion first, with its automatic strategy and the given time limit, and Vampire is asked only about the laws that E leaves undecided. On 500 such problems drawn from the questions of the fourth exploration of Section~\ref{sec:case3}, E decided 453 and Vampire 445, E decided 10 that Vampire did not and Vampire 2 that E did not, and the two never contradicted each other. When E decided, it did so within 0.34 seconds in 99\% of the cases and never took more than 1.04 seconds, whereas it used its whole limit when it did not; with 2 seconds for E first, the step cost half as much processor time as with both provers running in full, and no decision of the sample was lost. On the whole conclusions of the first step, E decided nothing that Vampire did not, so it is not used there. E's saturation sometimes succeeds where Vampire's does not: in the fourth exploration of Section~\ref{sec:case3}, it showed within a second that none of the 60 laws of the conclusion of an unconfirmed implication follows from its premise, which Vampire failed to show in 30 seconds per law even with its schedule for satisfiable problems. The row is then completed with the known implications, as the context completes its objects: the laws that the positive side implies are positive, and a law that would imply a law on the negative side is negative. The laws of the conclusion get the proving time limit (by default 10 seconds); the others, which decide nothing about the question but complete the row, get a separate \emph{row time limit} (option \code{-{}-vampire-row-time-limit}, by default the proving time limit). With row time limit~$0$, only the laws of the conclusion are asked about, and the other unknown laws stay unknown unless completion fills them in.

The row time limit matters when many questions end with a free magma that Vampire cannot fully decide: each of them then costs the row time limit for the laws outside the conclusion, which run in parallel on all cores. With the 149 laws of Section~\ref{sec:case3}, a premise that implies few of them leaves well over a hundred laws to ask about, so that every such question costs about the row time limit times the number of laws divided by the number of cores: about 3 minutes on 8 cores with the default 10 seconds, and under 20 seconds with a limit of 1 second.

A free magma is a powerful object: when its row is complete, its intent is the closure of $P$ in $\mathbb{K}_{\mathrm{Mag}}$, so it refutes every later question whose premise it satisfies and whose conclusion is not derivable. In the results, free magmas are named \code{F}$k$, the dual of $F_P$ being the free magma over the dual laws of $P$. The negative side of a free magma rests on Vampire's CounterSatisfiable verdicts. Those obtained from a finite model could be checked; those obtained by saturation cannot be checked by the library and are trusted.

\subsection{The completion expert}\label{sec:completion}

With \code{-{}-expert completion}, a question $P \to C$ is decided by Knuth--Bendix completion of the laws of $P$ into a convergent rewrite system $R$ (Section~\ref{sec:case31} describes how such a system certifies non-implications). If the two sides of every law of $C$ have the same normal form, the implication is accepted as proved. Otherwise the free magma over $P$ refutes it, and its row is complete: the normal forms decide every law. The system of each premise is checked as a certificate once, so that every negative entry is certified; positive entries rest on the soundness of completion. When completion fails, on an equation that the order cannot orient, such as commutativity, or at its limits of 400 rules and 10 seconds, the question goes to the Vampire expert if Vampire is installed, and is otherwise accepted as unconfirmed. The expert is meant for one-variable laws, for which completion has always succeeded; for laws in several variables it often fails. It has no pool of finite magmas, so \code{-{}-seed-magmas} does not apply to it.

\subsection{Seeding with known magmas}\label{sec:seeding}

Magmas found in earlier explorations (the \code{magmas.json} of saved results) can be reused in two ways, and so can the finite magmas of the ETP: files with one table such as \code{[[0, 1], [1, 0]]} per line, after the word \code{Table} or an equation number, as in the ETP's \path{equational_theories/Generated/All4x4Tables/data} and \path{data/smallest_magma_examples.txt}, are read as well (repeated tables once). They can be added to the expert's pool, which is checked before any solver runs. Alternatively, they can be added to the context before the exploration starts: they are evaluated on every law, and objects whose rows the others account for are dropped (an object is redundant if, for each law it violates, another remaining object satisfies everything it satisfies and violates that law). No question that the seeded objects refute is then asked.

%======================================================================
\section{Case study: sixteen classical laws}\label{sec:case16}
%======================================================================

The file \code{equations.json} lists the sixteen laws of Table~\ref{tab:laws16}. Some of them have two or more variables, so no background implications are derived; the only prior knowledge is the duality symmetry.

\begin{table}[ht]
\centering
\small
\begin{tabular}{@{}lll@{}}
\toprule
Law & Name & Equation \\
\midrule
$\Eq{2}$ & Singleton & $x = y$ \\
$\Eq{3}$ & Idempotence & $x = x * x$ \\
$\Eq{4}$ & Left-zero & $x = x * y$ \\
$\Eq{5}$ & Right-zero & $x = y * x$ \\
$\Eq{10}$ & Left-central identity & $x = x * (y * x)$ \\
$\Eq{16}$ & Steiner law 1 & $x * (x * y) = y$ \\
$\Eq{25}$ & Central identity & $x = (x * y) * x$ \\
$\Eq{26}$ & Steiner law 2 & $(y * x) * x = y$ \\
$\Eq{43}$ & Commutativity & $x * y = y * x$ \\
$\Eq{323}$ & Left-idempotent composition & $x * y = x * (x * y)$ \\
$\Eq{378}$ & Right-idempotent composition & $(y * x) * x = y * x$ \\
$\Eq{4396}$ & Left-alternative & $(x * x) * y = x * (x * y)$ \\
$\Eq{4435}$ & Flexible & $(x * y) * x = x * (y * x)$ \\
$\Eq{4473}$ & Right-alternative & $(y * x) * x = y * (x * x)$ \\
$\Eq{4512}$ & Associativity & $(x * y) * z = x * (y * z)$ \\
Medial & Medial (entropic) & $(x * y) * (z * w) = (x * z) * (y * w)$ \\
\bottomrule
\end{tabular}
\caption{The laws of \code{equations.json}. The names are those of the file.}\label{tab:laws16}
\end{table}

\subsection{Two explorations}

The laws were explored twice, once with each expert, with the default settings (Table~\ref{tab:runs16}).

\begin{table}[ht]
\centering
\small
\begin{tabular}{@{}lrr@{}}
\toprule
 & Z3, sizes 1--6 & Vampire, 10\,s / 1\,s / 10\,s \\
\midrule
Questions & 59 & 62 \\
Implications accepted & 27 & 27 \\
\quad of which unconfirmed & 2 & 0 \\
Implications found (with images) & 43 & 44 \\
Counterexamples & 32 & 35 \\
Free magmas & --- & 0 \\
Rejected implications & --- & 0 \\
Running time & 2 min 52 s & 2 min 51 s \\
\bottomrule
\end{tabular}
\caption{The two explorations of the sixteen laws. Of Vampire's 35 counterexamples, 27 are models found by Vampire and 8 come from the pool. The Vampire time limits are those for a proof, the quick model search and the per-law model search.}\label{tab:runs16}
\end{table}

The Vampire exploration proved every implication it accepted and refuted every other question with a checked finite magma, so its theory is exactly the set of implications between these laws that hold in all magmas. Its implications with a single premise agree with the ETP (Section~\ref{sec:etp}): fifteen of the sixteen laws, all but mediality, are ETP laws, and for each of the 210 ordered pairs of distinct ones, the theory contains the implication from the first to the second exactly when the ETP graph does. Its rules are listed in Table~\ref{tab:basis16}; each rule stands also for its dual, which swaps $\Eq{4}$ and $\Eq{5}$, $\Eq{10}$ and $\Eq{25}$, $\Eq{16}$ and $\Eq{26}$, $\Eq{323}$ and $\Eq{378}$, $\Eq{4396}$ and $\Eq{4473}$, and fixes the other laws.

\begin{table}[ht]
\centering
\small
\begin{tabular}{@{}rl@{\quad}rl@{}}
\toprule
1 & $\Eq{16}, \Eq{26} \to \Eq{43}$ &
15 & $\Eq{10} \to \Eq{378}$ \\
2 & $\Eq{4512} \to \Eq{4435}, \Eq{4473}$ &
16 & $\Eq{10}, \Eq{323} \to \Eq{4473}$ \\
3 & $\Eq{26}, \Eq{378} \to \Eq{4}$ &
17 & $\Eq{10}, \Eq{4435} \to \Eq{25}$ \\
4 & $\Eq{16}, \Eq{378} \to \Eq{10}, \text{Medial}$ &
18 & $\Eq{10}, \Eq{4473} \to \Eq{3}$ \\
5 & $\Eq{323}, \Eq{378}, \text{Medial} \to \Eq{4512}$ &
19 & $\Eq{10}, \Eq{4396} \to \Eq{323}$ \\
6 & $\Eq{43} \to \Eq{4435}$ &
20 & $\Eq{5} \to \Eq{4512}, \Eq{10}, \Eq{16}$ \\
7 & $\Eq{26}, \Eq{43} \to \Eq{16}$ &
21 & $\Eq{3}, \text{Medial} \to \Eq{4435}$ \\
8 & $\Eq{43}, \Eq{378} \to \Eq{323}$ &
22 & $\Eq{3}, \Eq{378} \to \Eq{4473}$ \\
9 & $\Eq{26}, \Eq{4473}, \text{Medial} \to \Eq{4512}$ &
23 & $\Eq{3}, \Eq{4473} \to \Eq{378}$ \\
10 & $\Eq{16}, \Eq{4473}, \text{Medial} \to \Eq{4512}$ &
24 & $\Eq{10}, \Eq{25} \to \Eq{4435}$ \\
11 & $\Eq{16}, \Eq{4435}, \Eq{4473} \to \Eq{4396}$ &
25 & $\Eq{25}, \Eq{4512} \to \text{Medial}$ \\
12 & $\Eq{43}, \Eq{4473} \to \Eq{4396}$ &
26 & $\Eq{25}, \Eq{43} \to \Eq{2}$ \\
13 & $\Eq{26}, \Eq{4512} \to \text{Medial}$ &
27 & $\Eq{2} \to$ everything \\
14 & $\Eq{43}, \Eq{4512} \to \text{Medial}$ & & \\
\bottomrule
\end{tabular}
\caption{The implications between the sixteen laws, as the rules printed by the program. Together with their duals, they entail every implication that holds in all magmas.}\label{tab:basis16}
\end{table}

\subsection{Where the experts differ}

The Z3 exploration accepted three implications that are false: one implication together with its dual, and one self-dual implication.
\begin{itemize}
  \item $\Eq{16}, \Eq{4435}, \Eq{4473} \to \Eq{4396}, \Eq{4512}, \text{Medial}$ (and its dual). Only $\Eq{4396}$ follows (rule~11). Associativity and mediality fail in \tocheck{a Steiner loop of order~10: a commutative loop in which every element is its own inverse, built from the Steiner triple system on nine points (the affine plane of order~3)}{the library only checked that the magma is a commutative loop with $x * x = e$ for its identity~$e$ and which laws it satisfies; the identification relies on the correspondence between Steiner loops of order $n + 1$ and Steiner triple systems on $n$ points, and on the uniqueness of the system on nine points.}. It satisfies $\Eq{16}$, $\Eq{26}$, $\Eq{43}$, $\Eq{4396}$, $\Eq{4435}$ and $\Eq{4473}$ but none of the other laws.
  \item $\Eq{3}, \Eq{16}, \Eq{26}, \Eq{43}, \Eq{4435} \to \text{Medial}$. It fails in \tocheck{the Steiner quasigroup of order~7, the idempotent commutative quasigroup of the Fano plane, in which $x * y$ is the third point on the line through $x$ and $y$}{the library only checked which laws the magma satisfies; the identification relies on idempotent totally symmetric quasigroups being Steiner quasigroups and on the uniqueness of the Steiner triple system on seven points.}.
\end{itemize}
Both counterexamples exceed the sizes the Z3 expert searches. Conversely, the two implications Z3 accepted as unconfirmed ($\Eq{323}, \Eq{378}, \text{Medial} \to \Eq{4396}, \Eq{4435}, \Eq{4473}, \Eq{4512}$, rule~5 before simplification, and $\Eq{43}, \Eq{4396}, \Eq{4435}, \Eq{4473}, \Eq{4512} \to \text{Medial}$) are true, and Vampire proved them. Every implication Vampire accepted is entailed by the Z3 theory, as it should be.

The first of the false implications also illustrates the per-law model search of step~6. Before that step was added, the Vampire expert left this question undecided: the free magma over the premise decided neither associativity nor mediality in time, and a model violating their conjunction was not found within the time limit. The model builder finds the size-10 Steiner loop in a few seconds when it looks for a violation of a single law, while the search for a violation of both needs more than a minute and a half.

%======================================================================
\section{Case study: one-variable laws of order at most four}\label{sec:case31}
%======================================================================

The file \code{equations\_1var\_order4.json} lists the 31 laws in one variable of order at most four, the trivial law $\Eq{1}$: $x = x$ included. Up to the order of the sides, they are all the equations $s = t$ between the 23 one-variable terms with at most four operations such that the two sides together have at most four operations: 22 laws $x = t$, seven laws $x * x = t$, the law $x * (x * x) = (x * x) * x$, and $\Eq{1}$. Since all laws have one variable, background implications can be derived.

\subsection{Background}

With order $N = 4$, transient order $N' = 10$ and premise size $K = 2$, the derivation closure yields 111 background implications. With premise size $0$, the canonical basis of the derivation closure has 95 implications; the two sets entail each other, so for these laws premises of size two lose nothing. The 111 implications have 108,635 closed sets.

\subsection{Results}

Table~\ref{tab:runs31} compares two Vampire explorations, which differ only in the order of the steps of the expert, with two explorations by the Z3 expert.

\begin{table}[ht]
\centering
\small
\begin{tabularx}{\textwidth}{@{}Xrrrr@{}}
\toprule
 & \multicolumn{2}{c}{Vampire} & \multicolumn{2}{c}{Z3, sizes 1--6} \\
\cmidrule(lr){2-3}\cmidrule(lr){4-5}
 & full search first & quick search first & no background & background \\
\midrule
Background implications & 111 & 111 & --- & 103 \\
Questions & 17,835 & 18,655 & 8,189 & 8,342 \\
Implications accepted & 0 & 0 & 887 & 842 \\
\quad of which false & 0 & 0 & 837 & 837 \\
Refuted by a finite model & 15,894 & 14,688 & 7,302 & 7,500 \\
Refuted by a free magma & 1,941 & 3,967 & --- & --- \\
Rejected or unconfirmed & 0 & 0 & 0 & 2 \\
Closed sets of the theory & 108,635 & 108,635 & 75,127 & 75,127 \\
Running time & 8 h 05 min & 2 h 25 min & not recorded & 3 h 18 min \\
\bottomrule
\end{tabularx}
\caption{Explorations of the 31 one-variable laws. ``Full search first'' runs a 10-second finite model search on the conjunction before computing the free magma; ``quick search first'' runs a 1-second search, then the free magma, then a 10-second search per law. ``Refuted by a finite model'' includes the 7 questions answered by the pool. The Z3 exploration with background used order 7 and no transient order. The Vampire running times are wall-clock times on an 8-core laptop; the second excludes a 40-minute sleep of the machine.}\label{tab:runs31}
\end{table}

\paragraph{No implication beyond the background.}
Both Vampire explorations refuted every question: no implication was accepted, and the background alone generates the theory. This is a mathematical result: \emph{every implication between one-variable laws of order at most four that holds in all magmas follows from the 111 background implications}, that is, it can be derived by equational reasoning through equations of order at most four, with single steps reaching order ten, from premises of at most two laws. Conversely, every non-implication is witnessed either by a finite magma, checked by the library, or by a free magma whose relevant negative entry Vampire found by saturation and Knuth--Bendix completion certifies (see below). The implications with a single premise agree with the ETP (Section~\ref{sec:etp}): all 31 laws are ETP laws, and for each of the 930 ordered pairs of distinct laws, the background implies the second from the first exactly when the ETP graph does. The final context of the second run has 37,177 objects (originals and duals) with pairwise distinct rows, and its 108,635 intents are exactly the closed sets of the background. Every row was decided completely: no free magma has an unknown entry.

\paragraph{Certifying the free magmas.}
The negative entries of the free magmas need not rest on Vampire. For laws $P$ and $s = t$, a rewrite system $R$ shows that $P$ does not imply $s = t$ if every rule of $R$ decreases in the Knuth--Bendix order, every critical pair of $R$ is joinable, every law of $P$ is joinable in $R$, and $s$ and $t$ have different normal forms: $R$ is then convergent, entails $P$, and does not entail $s = t$. These conditions are checked on $R$ itself, however it was found. The script \code{completion.py} finds $R$ by Huet's completion of $P$ with the Knuth--Bendix order of weight~1, and checks it. For the first Vampire exploration, completion terminated for every one of the 3,858 premises of its free magmas (originals and duals), within six seconds in all, and certified all 80,886 negative entries; the certificates are saved with its results (\path{free_magma_certificates.json}). It also certified all 167,042 negative entries of the 7,910 free magmas of the second exploration. As a check of the checker, none of the 38,712 positive entries of the free magmas of the first exploration is separated, and deciding by completion each of the 930 implications with a single premise between the 31 laws agrees with the ETP graph. With completion, the non-implications of this case study thus rest on checked finite magmas and checked rewrite systems, and no longer on Vampire.

\paragraph{Exploration by completion.}
The completion expert (Section~\ref{sec:completion}) decides every question of this case study by itself: two explorations with it, one with the background above and one without any background, never needed Vampire. With the background, it asked 54,575 questions, refuted all of them with free magmas and accepted nothing, in 8~min~15~s; without background, it asked 54,618 questions, refuted 54,568 with free magmas and proved 50 implications, 95 with their duals, in 8~min~57~s. Both theories entail every implication of the Vampire explorations and conversely, and have the same 108,635 closed sets, so this case study's result is confirmed by an expert that uses neither Vampire's proofs nor its saturation nor finite models. The explorations ask about three times as many questions as the Vampire ones, since the row of the free magma over $P$ is exactly the closure of $P$, so that it refutes only the question it answers and its dual, whereas a finite magma often refutes many; but each answer takes milliseconds.

\paragraph{The counterexamples.}
The finite counterexamples found in the first Vampire run have the following sizes (originals only):
\begin{center}
\small
\begin{tabular}{@{}lrrrrrrrrr@{}}
\toprule
Size & 2 & 3 & 4 & 5 & 6 & 7 & 8 & 9 & 10 \\
Magmas & 5 & 92 & 876 & 2,741 & 2,837 & 2,247 & 4,694 & 2,200 & 202 \\
\bottomrule
\end{tabular}
\end{center}
Most counterexamples are larger than six elements, and some implications have no finite counterexample that Vampire finds in time; for those the free magma is the counterexample. A typical example is question~685 of the second run, whose premise consists of $\Eq{1}$ and seven laws including $x = (((x*x)*x)*x)*x$ and $x*x = x*(x*(x*x))$: the free magma over the premise satisfies exactly the eight premise laws and refutes all 23 others at once.

\paragraph{What the Z3 explorations got wrong.}
The Z3 expert searches only up to six elements, and its explorations accepted 887 and 842 implications. Only 50 and 5 of them, respectively, follow from the background (the second exploration already had most valid implications as background); the remaining 837 are false in both cases. The false implications made the enumeration skip about a third of the closed sets, which is why the Z3 explorations asked fewer than half as many questions. In this domain, bounded finite-model search is a poor expert: an implication without small counterexamples is more often false than true.

\paragraph{Order of the steps.}
In the first Vampire exploration, a 10-second model search on the conjunction preceded the free magma. Most questions that ended with a free magma thus first spent 10 seconds in a hopeless search. Running a 1-second search first, the free magma next, and a 10-second search per law last cut the running time by a factor of more than three. The number of questions differs between the runs because different counterexamples were found, which changes the context and therefore which closed sets need a question; the theory is the same.

%======================================================================
\section{Case study: ETP laws of order at most two and three}\label{sec:case3}
%======================================================================

The ETP has decided every implication $E \to F$ between two of its laws, but an implication with a larger premise, such as $E_1, E_2 \to F$, is not determined by those: $E_1$ and $E_2$ may together imply a law that neither implies alone. Attribute exploration with the ETP graph as background (Section~\ref{sec:etp}) asks exactly about such premises. The two explorations of this section apply it to the ETP laws of order at most two and three.

Implications whose premise has several laws are among the future directions listed in the ETP paper~\cite[Section~14, item~(2)]{ETP25}, and were discussed on the ETP's Zulip channel from its start. There, Bruno Le Floch found with Prover9 that the conjunctions of the laws of order at most two fall into 59 equivalence classes, and computed all implications between them\footnote{\url{https://leanprover-community.github.io/archive/stream/458659-Equational/topic/Search.20space.20for.20(EqX.20.E2.88.A7.20EqY).20.E2.86.92.20EqZ.html}}. He has since determined all equivalence classes of conjunctions, and the implications between them, for 990 quasigroup laws~\cite{LF26a} and for semigroup laws of order up to four~\cite{LF26b}, and reports in the latter that the implication semilattice of conjunctions of magma laws of order up to three, with a few thousand equivalence classes, is almost determined in work in progress by other members of the ETP. The explorations of this section are therefore independent computations of known or nearly known results, by a different method: attribute exploration, which yields the canonical basis of the theory relative to the ETP background, and with the rules of order at most two proved in Lean.

\subsection{Choice of the laws}

Laws that imply each other are interchangeable in an implication, so it suffices to explore one law from each class of equivalent laws. The ETP lists 46 laws of order at most two, in 23 classes, and 410 laws of order at most three, in 146 classes. One law is chosen from each class: the law with the smallest number, or, if the class contains a law equal to its own dual, the first such law. If no law of a class is self-dual, the dual of the chosen law, which lies in the dual class (possibly the same class), is chosen as well and represents that class, so that the set is closed under duality. This gives 24 laws of order at most two (\code{equations\_order2.json}) and 149 of order at most three (\code{equations\_order3.json}). Both files were written by \code{etp.py}, which computes the classes from the ETP graph and lists the members of each representative's class in the file:
\begin{Verbatim}
python explorations/equational_theories/etp.py equations.txt \
    --max-order 3 --graph graph.json -o equations_order3.json
\end{Verbatim}
Each set includes $\Eq{1}$: $x = x$ and $\Eq{2}$: $x = y$, which hold in every magma and only in the trivial magma, respectively. The two explorations took 18 and 138 background implications from the ETP graph.

\subsection{Order at most two}

With the Vampire expert and the ETP background, the exploration of the 24 laws asked 45 questions. It found 19 counterexamples, 6 of them two-element magmas of the pool and 13 found by Vampire, of sizes 3 to 5, and proved 26 implications; nothing was left unconfirmed or rejected, and no free magma was needed. The theory has 59 closed sets, and each of its 26 rules has a premise of two laws. For example:
\begin{itemize}
  \item $\Eq{25}, \Eq{43} \to \Eq{2}$: a commutative magma satisfying $x = (x * y) * x$ is trivial;
  \item $\Eq{26}, \Eq{29} \to \Eq{43}$: the laws $x = (x * y) * y$ and $x = (y * x) * y$ together imply commutativity;
  \item $\Eq{25}, \Eq{8} \to \Eq{3}$: the laws $x = (x * y) * x$ and $x = x * (x * x)$ together imply idempotence.
\end{itemize}
Since every question was decided, this is the complete implicational theory of the ETP laws of order at most two: every implication between them that holds in all magmas follows from the 26 rules, the 18 single-premise implications of the ETP graph, duality, and the equivalence of the laws of each class. Its 59 closed sets agree with Le Floch's 59 equivalence classes of conjunctions.

\subsection{Order at most three}

The 149 laws have far more closed sets, and many premises combine laws whose consequences Vampire finds hard to decide. Four explorations were run, all with the ETP background, the context seeded with the ETP's finite magmas (Section~\ref{sec:etp}; 233 of 1374 tables were kept), a row time limit of 1 second, and otherwise the default limits. The second also used Mace4 in the per-law model search, for 30 seconds per law, the third, in addition, the ETP's refutations as objects, and the fourth, in addition, asked about the duals of unconfirmed implications instead of adding them (Section~\ref{sec:symmetry}). Table~\ref{tab:runs3} compares them. With the ETP data fetched into \code{etp} (Section~\ref{sec:etp}), the fourth was run as
\begin{Verbatim}
python explorations/equational_theories/explore.py \
    --equations explorations/equational_theories/equations_order3.json \
    --expert vampire --vampire-row-time-limit 1 --mace4-time-limit 30 \
    --etp-graph etp/graph.json --etp-equations etp/equations.txt \
    --etp-refutations --ask-unconfirmed-duals --seed-context etp/magmas/*.txt
\end{Verbatim}


\begin{table}[ht]
\centering
\small
\begin{tabularx}{\textwidth}{@{}Xrrrr@{}}
\toprule
 & Vampire & and Mace4 & and refutations & and duals asked \\
\midrule
Questions & 1,310 & 1,326 & 1,337 & 1,376 \\
\quad proved & 736 & 746 & 746 & 748 \\
\quad refuted by a finite model & 437 & 440 & 444 & 444 \\
\quad refuted by a free magma & 75 & 76 & 72 & 74 \\
\quad unconfirmed & 57 & 59 & 70 & 105 \\
\quad rejected & 5 & 5 & 5 & 5 \\
Accepted implications & 807 & 802 & 824 & 860 \\
\quad confirmed & 750 & 754 & 765 & 767 \\
\quad unconfirmed & 57 & 48 & 59 & 93 \\
\qquad of which known to be false & 6 or more & 3 & 1 & 1 \\
Canonical basis relative to the background, up to duality & 793 & 795 & 806 & 806 \\
\quad of the confirmed implications alone & 750 & 754 & 765 & 767 \\
Closed sets of the theory & 3,650 & 3,663 & 3,667 & 3,671 \\
\quad without the unconfirmed implications & 29,415 & 28,927 & 11,107 & 11,097 \\
Rounds & 1 & 2 & 2 & 2 \\
Running time & 2 h 48 min & 4 h 34 min & 4 h 42 min & 6 h 19 min \\
\bottomrule
\end{tabularx}
\caption{Explorations of the 149 representatives of the ETP laws of order at most three, with the ETP graph as background and the ETP's finite magmas in the context: by the Vampire expert, with Mace4 in the per-law search as well, with the ETP's refutations in the context as well, and asking about the duals of unconfirmed implications as well. The questions are counted by answer, so that a question asked again in a later round counts twice. The accepted implications are those the expert accepted, without the background and without the duals that the symmetry adds; in the fourth exploration, the duals of unconfirmed implications are asked about and accepted in their own right, so that its 93 unconfirmed implications are 60 up to duality. Those known to be false were shown false by later explorations or tests, and they remain in their exploration's result. The canonical basis is that of the theory relative to the background, counted up to duality: in all, it has 1,552, 1,557, 1,579 and 1,579 implications, of which 33 or 34 have self-dual premises. The second line gives the canonical basis, relative to the background and up to duality, of the theory that the confirmed implications and their duals generate. In the first three explorations, the rules printed by the program (Section~\ref{sec:details}) are as many as the canonical basis up to duality; in the fourth, which does not map unconfirmed implications by duality, they are 842, its 93 unconfirmed implications being 60 up to duality. The closed sets without the unconfirmed implications are counted without their duals as well, but with the duals of the confirmed implications, a few of which the fourth exploration did not record because unconfirmed implications entailed them. The running times are wall-clock times on an 8-core laptop; the second excludes a sleep of the machine of about an hour.}\label{tab:runs3}
\end{table}

\paragraph{Where Vampire gets stuck.}
Most hard questions involve $\Eq{255}$: $x = ((x * x) * x) * x$ or its dual $\Eq{47}$: $x = x * (x * (x * x))$, combined with laws in two or more variables. In the second exploration, 44 of the 48 unconfirmed implications and all five rejected implications have one of these two laws in their premise. Their premises are often large: the rejected implications have 16 to 18 laws in their premise.

\paragraph{Testing the unconfirmed implications.}
After the first exploration, each law of the conclusion of each of its 19 unconfirmed questions (109 laws in all) was tried again, separately, with longer limits:
\begin{itemize}
  \item Vampire, 60 seconds in portfolio mode and 60 seconds with its model builder, first on the premise alone and then with the one-variable laws that the premise derives (Section~\ref{sec:background}) added as lemmas: none of the 109 laws was decided;
  \item Prover9, 60 seconds: no proof;
  \item Mace4, 60 seconds: models of size 16 violating 34 of the laws, found within 0.2 to 31 seconds. These settled six of the 19 questions as false, all their unproved laws being violated, and three magmas sufficed for all six.
\end{itemize}
So at least six of the 57 unconfirmed implications of the first exploration are false, and the enumeration, which they steered, skipped closed sets. Its 3,650 closed sets are therefore not those of the true theory.

\paragraph{The second exploration.}
With Mace4 in the per-law search, the expert found the counterexamples of size 16 during the exploration itself, together with models of sizes 11 and 13 that Vampire had missed (17 tables in all, counting duals). One of them refuted an unconfirmed implication accepted 260 questions earlier; the exploration resumed at that question (Section~\ref{sec:rounds}), which made it run in two rounds. In the end, 48 unconfirmed implications and 5 rejected implications remain, and 46 of the 48 have the same premise as an unconfirmed implication of the first exploration. No counterexample to them is known, but they are not proved either.

\paragraph{Infinite counterexamples.}
Three of these 48 unconfirmed implications are nevertheless false:
\begin{itemize}
  \item $\Eq{1}, \Eq{255}, \Eq{261} \to \Eq{307}$;
  \item $\Eq{1}, \Eq{255}, \Eq{274} \to \Eq{47}, \Eq{99}, \Eq{118}, \Eq{203}, \Eq{222}$;
  \item $\Eq{1}, \Eq{47}, \Eq{99}, \Eq{203}, \Eq{229}, \Eq{255} \to \Eq{73}, \Eq{125}$.
\end{itemize}
In each, one law of the premise ($\Eq{261}$, $\Eq{274}$ and $\Eq{229}$) implies the others, and the ETP refutes that it implies some law of the conclusion: $\Eq{261} \not\to \Eq{307}$, $\Eq{274} \not\to \Eq{47}$, and $\Eq{229} \not\to \Eq{73}$. The ETP's counterexamples are infinite: for the duals $\Eq{65}$ (the Asterix law of the ETP), $\Eq{73}$ and $\Eq{118}$ of these laws, it builds them greedily on $\mathbb{N}$ or on a free group. For the first and the third implication, no finite counterexample exists, since the ETP proves in Lean that the duals of their single-premise refutations hold in finite magmas (\path{Finite.Equation65_implies_Equation359}, \path{Finite.Equation118_implies_Equation274} and \path{Finite.Equation118_implies_Equation222}); Mace4 also finds no magma of size at most 40 that satisfies $\Eq{261}$ and violates $\Eq{307}$. Whether the second holds in finite magmas was not checked. No finite search could have refuted the first and the third, and since they steered the enumeration, the theory of the second exploration is not the theory of all magmas even if its other unconfirmed implications are true.

\paragraph{The third exploration.}
The third exploration added the ETP's refutations to the context (Section~\ref{sec:etp}). Of the 20,916 implications between two of the 149 laws that the ETP refutes, 10 remained after the reduction against the 233 magmas, five and their duals: $\Eq{261} \not\to \Eq{307}$, $\Eq{274} \not\to \Eq{47}, \Eq{99}, \Eq{203}$ and $\Eq{229} \not\to \Eq{73}$, exactly those behind the three false implications. They were not accepted. Instead, the exploration asked about the larger premises that they had covered, and accepted seven implications, counting duals, that the second theory does not entail, all unconfirmed, such as $\Eq{1}, \Eq{255}, \Eq{274}, \Eq{280} \to \Eq{43}, \Eq{73}, \Eq{75}, \Eq{125}, \Eq{229}, \Eq{264}$. Apart from these and the three false implications with their duals, each theory entails every implication of the other. The number of unconfirmed implications grew from 48 to 59, since each false implication stood for several finer questions that the expert could not decide either.

\paragraph{The fourth exploration.}
In the first three explorations, the dual of each unconfirmed implication was added with it and steered the enumeration as well. The fourth asked about these duals instead (Section~\ref{sec:symmetry}). It asked every question of the third exploration and 38 more, of which 35 were left unconfirmed, mostly the duals themselves, so that it took 1 h 37 min longer. Its answers to the common questions were the same but for one: question 262, $\Eq{1}, \Eq{255}, \Eq{257}, \Eq{263}, \Eq{307}, \Eq{310} \to \Eq{260}, \Eq{266}, \Eq{323}$, unconfirmed in the third exploration, was refuted by a magma of size 13 that Vampire found within its 10 seconds this time. The third theory therefore contains a false implication and its dual; the fourth proves the rest of that question, $\Eq{323}$ added to the premise, and confirms one large implication more. Apart from these, the two theories are the same, and the fourth found no unproved implication that the third theory lacks. Asking about the duals thus made the result slightly more complete in this case, but mainly by the difference at question 262 rather than by easier questions.

\paragraph{Posing the questions differently.}
The unconfirmed questions of the fourth exploration were also put to the provers in other forms, with the expert's time limits, to see whether the form of a question matters:
\begin{itemize}
  \item with minimal premises, the fewest laws whose closure under the background and confirmed implications is the premise (2.4 laws instead of 13.6 on average): no proof among the 1,652 unproved laws, and no counterexample among 222 of them sampled; on 444 refuted questions, the full premises found three large models (sizes 15 and 16) that the minimal ones missed;
  \item with premises relaxed to the laws not added by unconfirmed implications: only one question had such laws, and it was not refuted;
  \item with the other members of the ETP equivalence class of each unproved law as the goal (1,865 goals for 294 laws), and the classes of the premise laws added to the premise: no proof and no counterexample.
\end{itemize}
Longer runs on single questions did not help either: the conjecture that, given $\Eq{3}$ and $\Eq{263}$, $\Eq{385}$ is equivalent to $\Eq{117} \wedge \Eq{159}$, of which one half is proved, resisted Vampire for 10 minutes, Twee for 30 minutes and Mace4 up to sizes 15 and 16.

\paragraph{What the results say.}
In the fourth exploration, the 767 confirmed implications, together with the background and their duals, are proved. Every non-implication witnessed by a finite magma is checked, those witnessed by free magmas rest on Vampire's saturation, and those witnessed by the ETP's refutations rest on the ETP. The 3,671 closed sets of the theory would be correct if the 93 unconfirmed implications were true, but two of them, an implication and its dual, are false: the E prover showed that none of the 60 laws of their conclusions follows from their premises (Section~\ref{sec:free}). E decided none of the laws of the other 91 within 10 seconds each, and showed for each of the 5 rejected implications that none of their conclusion laws follows, so that free magmas refute them. The confirmed implications, the background and their duals generate a theory with 11,097 closed sets, an upper bound on the number of closed sets of the true theory, since each confirmed implication holds. None of the 93 is refuted by a single-premise refutation of the ETP, which are all in the context, and whether they hold in finite magmas is not known. Settling the unconfirmed implications and finding counterexamples to the rejected ones needs either proofs that Vampire and Prover9 do not find within a minute, or counterexamples that Vampire and Mace4 do not find: finite magmas beyond their reach, or free magmas whose rows Vampire cannot decide. A comparison with the semilattice of the work in progress mentioned in~\cite{LF26b} would settle many of them, or reveal errors on either side.

%======================================================================
\section{Case study: semigroup laws of order at most four}\label{sec:semigroups}

Le Floch~\cite{LF26b} determined all implications between conjunctions of the laws of order at most four on semigroups: the 653 ETP laws whose two sides are right-associated, which are the ETP's forms of the laws on an associative operation, fall into 122 classes of equivalent laws and generate 456 theories (``associative theories''), found with Prover9, Mace4 and Vampire, with the proofs formalized in Lean. This section reproduces that result by attribute exploration, from the ETP's data alone.

\subsection{Setup}

\paragraph{Objects.} Option \code{-{}-axiom 4512} makes every object a semigroup: associativity is concluded from the empty premise in the background, so that every question's premise contains it, and of the ETP's 1374 finite magmas only the 24 associative ones are kept, of which the 12 that are not redundant are seeded into the context. The duality compares laws as words (Section~\ref{sec:duality}).

\paragraph{Laws.} The 653 laws were selected from the ETP's list by the syntactic condition alone, both sides right-associated and order at most four; their number agrees with Le Floch's. The ETP graph, whose implications hold in all magmas and so in all semigroups, splits them into 219 classes. Its equivalences are only those that hold without associativity; the ones that associativity adds, which lead to Le Floch's 122 classes, are left to the exploration. One law of each class is enough, but the reversal of a law's words, its dual among semigroups, is equivalent to its mirror image only modulo associativity and need not lie in the class of that image, so the 219 first laws of the classes are not closed under duality. Adding their reversals (28 laws) and associativity gives the 248 laws of \code{equations\_semigroup.json}, written by \code{etp.py -{}-max-order 4 -{}-semigroup} (Section~\ref{sec:guide}). The ETP graph gives 228 background implications between them.

\paragraph{Expert.} The Vampire expert ran with the settings of \code{configs/semigroup.toml}: 3 seconds on all cores for the first proof attempt, E for 2 seconds before Vampire in the free magma's row, row time limit~$0$, and, in the per-law model search (Section~\ref{sec:vampire}), 300 seconds for Vampire's model builder per law and 600 seconds for Mace4's size sweep, both with simplified premises:
\begin{Verbatim}
python explorations/equational_theories/explore.py \
    --config explorations/equational_theories/configs/semigroup.toml
\end{Verbatim}

\subsection{Results}

\begin{table}[ht]
\centering
\small
\begin{tabularx}{\textwidth}{@{}Xr@{}}
\toprule
Laws explored (of the 653, with associativity) & 248 \\
Background implications from the ETP graph & 228 \\
Questions & 251 \\
\quad proved & 204 \\
\quad refuted by a finite model found for the question & 40 \\
\quad refuted by a model in the pool & 7 \\
\quad refuted by a free magma, unconfirmed or rejected & 0 \\
Canonical basis relative to the background and duality & 204 \\
\quad relative to the ETP graph's implications alone, up to duality & 214 \\
\quad relative to no background, up to duality & 308 \\
Closed sets of the theory & 456 \\
Rounds & 1 \\
Running time & 18 min \\
\bottomrule
\end{tabularx}
\caption{The exploration of the semigroup laws of order at most four. The 47 counterexamples are semigroups of sizes 3 to 12, 16, 27 and 32, found by Vampire's model builder (30) and Mace4 (17); seven of them, found while the search went on for other laws of a question, answered a later question from the pool. Duality acts on the background as well: the dual of an implication of the ETP graph between right-associated laws reverses their words, holds only in semigroups and is not in the graph, and the exploration adds these duals, 86 implications, to the background it works with. Relative to that background, and with every question decided in one round, the canonical basis up to duality consists of the 204 accepted implications (336 implications in all, 72 of them self-dual). Relative to the graph's implications and associativity alone, it has 214 up to duality (343 in all), and relative to no background, for the whole theory, 308 (508 in all). The program prints 204 rules. The running time is the wall-clock time on an 8-core laptop.}\label{tab:semigroups}
\end{table}

Every question was decided (Table~\ref{tab:semigroups}): the 204 accepted implications are proved, and the 47 counterexamples are finite semigroups, checked against every law, so that no answer rests on a prover's saturation.

\paragraph{Comparison with Le Floch's result.} A law of the 653 holds in a closed set of the exploration when the first law of its ETP class does. So counted, the 456 closed sets are exactly Le Floch's 456 theories, and the classes of laws with the same closure are his 122, except that the exploration puts $x = x$ and associativity in one class, associativity being assumed. The exploration thus reproduces the classes, the theories and the implications between them without using any of his data.

\paragraph{The hardest questions.} Four questions needed large models. Two have the premises of theories that Le Floch singles out~\cite[Section~4]{LF26b}: the premise of question~216 generates the theory of $\Eq{56} \wedge \Eq{75}$, whose smallest non-commutative model has 27 elements, and that of question~159 the theory of $\Eq{440} \wedge \Eq{513}$, whose smallest model violating $\Eq{504}$, $x = y(y(x(yy)))$, has 32. Mace4's sweep found both models. Questions~208 and~233 needed models of 16 elements, found by Mace4 and by Vampire's model builder.

\subsection{How the model search was extended}

The per-law model search of Section~\ref{sec:vampire} took its present form during this case study. A first exploration of Le Floch's 122 class representatives with associativity, without duality and with the former per-law search (10 seconds for Vampire's model builder and 30 for Mace4, both trying the sizes in increasing order, on the full premise), accepted three false implications as unconfirmed, which, written as words,
\begin{itemize}
  \item $xx = xyzy$, $xyz = xzy \to xx = xyzw$;
  \item $xx = yyxx$, $xyy = yxx$, $xyz = zyx \to xx = yxzw$;
  \item $xx = yzyx$, $xyz = yzx \to xx = yxzw$,
\end{itemize}
for which the premises are given in simplified form. They hid 8 of the 456 theories, and two true implications were never asked. One semigroup of 16 elements refutes all three; for the third premise, Mace4 shows that no smaller one exists. Separate tests showed what finds such models:
\begin{itemize}
  \item \emph{A fixed size.} Mace4 searching sizes in increasing order reached size~13 in 15 minutes without a model, while at size~16 alone it found one in 213 seconds. Ruling out the sizes below the smallest model is what takes the time, roughly doubling from one size to the next.
  \item \emph{The simplified premise.} At size~16, Mace4 needed 213 seconds with the simplified premise of the first implication and 897 with its full premise of 25 laws; for the third implication, Vampire's model builder found a model in 112 seconds with the simplified premise and none in 15 minutes with the full one, which it left at size~8. For magmas of order three, the full premise had done slightly better (Section~\ref{sec:case3}), hence the option.
  \item \emph{Not saturation.} E and Vampire, 10 minutes each, saturated none of the three problems, so longer proof attempts in the free magma's row would not have helped.
\end{itemize}
Larger sizes are not hopeless for laws whose two sides are products, such as those of the three implications: adding a copy of an element, with the same row and column, preserves such a law and its failures, so that a counterexample of size $n$ gives counterexamples of every larger size, and a search at any size at least the smallest succeeds. The size sweep exploits this on all cores without guessing the size, and the search goes on for the laws a model does not violate, so that the pool answers the next questions about the same premise. With these changes and duality, an exploration of Le Floch's representatives together with the reversals of 11 of them, which his classes contain, already gave his 456 theories; the exploration of this section chooses its laws from the ETP's classes instead.

%======================================================================
\section{User's guide}\label{sec:guide}
%======================================================================

\subsection{Installation}

The library needs Python 3.11 or later. In a clone of the repository:
\begin{Verbatim}
pip install -e ".[examples]"     # the library, Z3 and Jupyter
\end{Verbatim}
The Vampire expert additionally needs the \code{vampire} executable on the path, for example via Homebrew (\code{brew install vampire}) or a release from \url{https://vprover.github.io}. The explorations described here used Vampire 5.1.0. The optional Mace4 search needs the \code{mace4} executable, which comes with Prover9 (\code{brew install prover9}).

\subsection{Describing the laws}

Write the laws to a JSON file as in Section~\ref{sec:et}. Every entry needs an \code{equation} string; \code{name} and \code{id} are optional, and laws without an id are identified by their name or their text in the results. Terms use variables (single letters or identifiers), \code{*} and parentheses; an unparenthesized product associates to the left, so \code{x * y * z} is read as $(x * y) * z$. The set must be closed under duality: for every law, a law equal to its dual up to renaming variables and swapping sides must be in the file, or the program stops with an error naming the missing dual. Derived background implications (Section~\ref{sec:background}) require all laws to have the same single variable; background from the ETP (Section~\ref{sec:etp}) does not. To explore ETP laws chosen by number, let \code{etp.py} write the file:
\begin{Verbatim}
python explorations/equational_theories/etp.py equations.txt 3 43 4512 \
    -o my_laws.json
\end{Verbatim}
where \code{equations.txt} is the ETP's list of laws; the duals of the chosen laws are added. With \code{-{}-max-order K -{}-graph graph.json} instead of numbers, it writes one law of each class of equivalent ETP laws of order at most $K$, closed under duality (Section~\ref{sec:case3}), each with the members of its class. With \code{-{}-semigroup} in addition, it chooses among the right-associated laws, the ETP's forms of laws on semigroups: the first law of each class, the reversals of these laws, so that the laws are closed under duality among semigroups, and associativity; the file \code{equations\_semigroup.json} lists the 248 laws so chosen for order at most four, to be explored with \code{-{}-axiom 4512} (preset \code{configs/semigroup.toml}).

\subsection{Running an exploration}

From the root of the repository:
\begin{Verbatim}
python explorations/equational_theories/explore.py \
    --equations my_laws.json --expert vampire
\end{Verbatim}
Table~\ref{tab:options} lists the options. They can also be read from a TOML file with \code{-{}-config FILE}, keyed by option name without the dashes (\code{eprover-time-limit = 2}), the options given on the command line overriding the file. Paths are relative to the directory \code{explore.py} is run from, and patterns among the seed files, such as \code{etp/magmas/*.txt}, are expanded. The file \path{explorations/equational_theories/configs/order3.toml} holds the settings for the laws of order at most three, with the ETP data fetched into \code{etp} (Section~\ref{sec:etp}). Every run saves all its options as \code{settings.toml} in its results directory, so that \code{-{}-config} with that file runs it again.

\begin{table}[ht]
\centering
\small
\begin{tabularx}{\textwidth}{@{}>{\ttfamily}l l X@{}}
\toprule
\normalfont Option & Default & Meaning \\
\midrule
-{}-equations PATH & \code{equations.json} & the laws to explore \\
-{}-expert z3|vampire|completion & \code{z3} & which expert answers the questions \\
-{}-axiom LAW... & & laws every object satisfies, such as associativity \\
-{}-no-duality & off & do not use the duality symmetry \\
-{}-vampire-prove-time-limit S & 10 & seconds for a proof, and for each law of the conclusion in a free magma's row \\
-{}-vampire-first-time-limit S & prove limit & seconds for the first proof attempt, on all cores; 3 suffices \\
-{}-vampire-row-time-limit S & prove limit & seconds for each other law of a free magma's row; $0$ leaves them unknown \\
-{}-vampire-quick-model-time-limit S & 1 & seconds for the first finite model search \\
-{}-vampire-model-time-limit S & 10 & seconds for each per-law finite model search \\
-{}-mace4-time-limit S & 0 & seconds for Mace4's size sweep in the per-law search, in all; $0$ leaves Mace4 out \\
-{}-model-premise full|simplified|both & \code{full} & the premise the per-law model searches use \\
-{}-eprover-time-limit S & 0 & seconds for E first on each law of the conclusion in a free magma's row, Vampire then getting the rest; $0$ leaves E out, 2 suffices \\
-{}-min-search-size N & 1 & smallest table size Z3 tries \\
-{}-max-search-size N & 6 & largest table size Z3 tries \\
-{}-z3-timeout-ms MS & 30000 & Z3 time per size; $0$ means unbounded \\
-{}-background-order N & off & derive background implications through equations of order at most $N$ (at least the largest order of a law) \\
-{}-background-transient-order M & $N$ & let single steps reach order $M$ \\
-{}-background-premise-size K & 2 & largest premise of a background implication; $0$ computes the canonical basis \\
-{}-etp-graph PATH & & the ETP's implication graph, as background; needs \code{-{}-etp-equations} \\
-{}-etp-equations PATH & & the ETP's list of laws, to match the explored laws \\
-{}-etp-refutations & off & add the graph's refutations to the context; needs \code{-{}-etp-graph} \\
-{}-ask-unconfirmed-duals & off & ask about the duals of unconfirmed implications instead of adding them \\
-{}-seed-magmas PATH... & & add magmas to the expert's pool \\
-{}-seed-context PATH... & & add magmas to the context before starting \\
-{}-report-every N & 1 & print every $N$th question; $0$ prints none \\
-{}-results-dir PATH & new directory & where to save the results \\
-{}-config FILE & & read options from a TOML file, the command line overriding it \\
\bottomrule
\end{tabularx}
\caption{Options of \code{explore.py}. The Z3 options apply to the Z3 expert only, and the Vampire options to the Vampire expert only.}\label{tab:options}
\end{table}

\subsection{Choosing the expert and the time limits}

\begin{itemize}
  \item Use the \textbf{Vampire expert} when the result should be correct. Its confirmed implications are proved, and its counterexamples are checked magmas or free magmas. If it accepts something as unconfirmed or rejects an implication without a counterexample, the summary says so.
  \item Use the \textbf{completion expert} for one-variable laws: it decides every question it can complete, in minutes rather than hours, and certifies its counterexamples.
  \item Use the \textbf{Z3 expert} for a quick first impression, or when Vampire is not available. Treat its accepted implications as conjectures: as the case studies show, they may be false.
  \item Longer \textbf{proving time limits} help when implications are true but hard to prove. The laws of the conclusion in a free magma's row get the same limit, and a law Vampire cannot decide in time stays unknown.
  \item A short \textbf{row time limit}, such as 1 second, speeds up explorations with many laws, where most of a free magma's row is outside the conclusion. The price is more unknown entries, so a free magma refutes fewer later questions by itself.
  \item Longer \textbf{model time limits} help when counterexamples are large. The per-law search usually finds a model faster than the search on the whole conclusion. Adding Mace4 with \code{-{}-mace4-time-limit} costs up to that many seconds for every question that reaches the per-law search, which is worth it when Vampire's model builder misses counterexamples of moderate size, such as the magmas of size~16 that refuted several unconfirmed implications between laws of order at most three.
  \item Vampire runs in parallel on all cores when it computes a free magma's row or searches for models law by law, so other work on the machine slows the exploration down.
\end{itemize}

\subsection{Background knowledge}

For one-variable laws, background implications usually pay off: they are computed in seconds and remove the questions whose answer is a derivation. Start with \code{-{}-background-order} equal to the largest order of the laws and a transient order of about twice that; the 31-law exploration used 4 and 10. A larger order finds more implications at a cost that grows quickly. A premise size of 2 is usually enough; compare with \code{-{}-background-premise-size 0} on a small set of laws if in doubt.

For laws the ETP lists, pass its implication graph and its list of laws:
\begin{Verbatim}
python explorations/equational_theories/explore.py \
    --equations my_laws.json --expert vampire \
    --etp-graph graph.json --etp-equations equations.txt
\end{Verbatim}
This background costs nothing to compute and leaves only questions whose premise has several laws. It helps most when the implications between pairs of laws are hard to settle. Add \code{-{}-etp-refutations} to put the graph's refutations into the context as well (Section~\ref{sec:etp}); with the ETP's finite magmas seeded, only refutations that they do not account for are kept, such as those by infinite magmas.

\subsection{Reusing magmas}

To start a new exploration with the magmas of an earlier one, pass its \code{magmas.json} to \code{-{}-seed-context} (fewer questions) or to \code{-{}-seed-magmas} (the same questions, answered faster). Free magmas are not reused. The ETP's files of tables can be passed in the same way. Seeded into the context of the exploration of Section~\ref{sec:case3}, the 1374 distinct tables of \path{All4x4Tables} and \path{smallest_magma_examples.txt} were reduced to 233 magmas.

\subsection{Reading the output}

The program prints every question with its answer and writes the same text to \code{exploration.log}:
\begin{Verbatim}
Number of implications: 111
Number of objects: 1358
Question 685: Eq1: x = x, Eq1629: ... -> Eq8: x = (x * (x * x)), ...
Counterexample found: Free magma over {Eq1, Eq1629, Eq1832, ...}
Satisfies: ...
Does not satisfy: ...
\end{Verbatim}
The counts are those before the answer; the objects include the duals. An accepted question is followed by \code{Confirmed} or \code{Unconfirmed}, and an undecided one by \code{Rejected} or \code{Unconfirmed} together with the laws it proved. At the end, the program prints a summary (questions, accepted implications, counterexamples, rejected and conflicting implications, rounds, and the Vampire expert's answers by kind), the rules, any rejected, conflicting and refuted unconfirmed implications, and a few counterexamples.

The results directory contains:
\begin{description}
  \item[\code{settings.toml}] every option of the run, with the commits of the library and of the ETP data, for \code{-{}-config};
  \item[\code{exploration.log}] everything printed;
  \item[\code{implications.txt}, \code{implications.json}] every implication found, together with its source, which is one of \code{background}, \code{confirmed}, \code{unconfirmed}, \code{mapped} and \code{mapped unconfirmed}; then the refuted unconfirmed implications, the rejected implications, and the conflicting ones with the rejected implication each contradicts;
  \item[\code{rules.txt}] everything found, simplified into rules, each with its provenance: the background, simplified the same way, and the rules printed at the end, each rule followed by its dual (\code{mapped background}, \code{mapped} or \code{mapped unconfirmed}, from the rule it is the image of) unless it is its own; together they entail every implication found. Each law is written out by its number and text, or by its two words when associativity is assumed;
  \item[\code{magmas.json}] the finite magmas of the context with their tables, under the names \code{M}$k$ and \code{M}$k$\code{\^{}op};
  \item[\code{free\_magmas.json}] the free magmas \code{F}$k$, each with its premise and its known positive and negative laws;
  \item[\code{witnesses.json}] the witnesses \code{W}$k$ of implications with one attribute in their conclusion that the expert showed false without a counterexample, each with the implication it refutes and its row;
  \item[\code{etp\_refutations.json}] with \code{-{}-etp-refutations}, the objects \code{R}$k$ standing for refutations of the ETP graph, each with the implication it refutes and its known positive and negative laws;
  \item[\code{context.cxt}] the objects with complete rows, in Burmeister's format;
  \item[\code{partial\_context.cxt}] all objects in the same format, with \code{?} for unknown entries.
\end{description}
The program saves the results only at the end. If it is interrupted, only the log remains.

\subsection{Using the library from Python}

The runner is a thin layer over the library, which can be used directly:
\begin{Verbatim}
from conceptual_exploration import explore_in_rounds, report_every, simplified_rules
from explorations.equational_theories.magma import ETP
from explorations.equational_theories.vampire import VampireExpert
from explorations.equational_theories.background import background_implications

laws = ETP.load_equations("my_laws.json")
expert = VampireExpert(attributes=laws, prove_time_limit=10,
                       quick_model_time_limit=1, model_time_limit=10)
result = explore_in_rounds(
    laws, expert,
    background_implications=background_implications(laws, 4, 2, 10),
    mappings=[ETP.get_duality_mapping(laws)],
    evaluate_all=True,
    on_question=report_every(100),
)
for rule in simplified_rules(result.base):
    print(rule)
print(result.base.open_questions, expert.answers)
\end{Verbatim}
\code{run\_magma\_exploration} in \code{explore.py} does the same with the options above and saves the results; the notebook \code{explore.ipynb} shows smaller examples. Another expert can be plugged in by subclassing \code{Expert} and implementing \code{validate}, which returns a refuting \code{PartialObject}, \code{None} to accept, or an \code{Undecided} value.

\subsection{Practical remarks}

\begin{itemize}
  \item \textbf{Size.} The number of questions grows with the number of closed sets of the final theory, which may be large: 108,635 for 31 laws without a single accepted implication. Expect hours for such an exploration and minutes for one with a few dozen closed sets per law.
  \item \textbf{Long runs.} Keep the terminal open. Putting the machine to sleep is harmless in principle, but a Vampire run in progress at that moment may time out on waking, since the time limits are wall-clock times; this can at worst turn a refutation into an undecided answer, which the summary reports.
  \item \textbf{Reproducibility.} Equations are sorted before they are given to the solvers, so a Z3 exploration is reproducible. Vampire's time limits make a Vampire exploration depend on the machine's speed: a different run may find different counterexamples and ask a different number of questions, but the theory is the same unless a question is left undecided.
  \item \textbf{Machine load.} For the same reason, other work on the machine can turn a refutation into an unconfirmed answer: in an order-3 exploration run alongside other provers, a counterexample of size 15 that Vampire had found in three earlier explorations was missed, and the false implication was refuted only 22 questions later. Run explorations on an otherwise idle machine.
  \item \textbf{Trust.} Finite counterexamples are checked by the library. Proofs and saturation-based CounterSatisfiable verdicts are Vampire's; independent confirmation requires checking Vampire's proofs, as \code{lean\_proofs.py} does for the rules of order at most two, or the free magmas' negative entries, as \code{completion.py} does when Knuth--Bendix completion of their premises succeeds.
\end{itemize}

\begin{thebibliography}{ETP25}
\bibitem[ETP25]{ETP25} Equational Theories Project contributors. \emph{The Equational Theories Project: Advancing collaborative mathematical research at scale}. arXiv:2512.07087, 2025. \url{https://arxiv.org/abs/2512.07087}
\bibitem[GO16]{GO16} Bernhard Ganter and Sergei Obiedkov. \emph{Conceptual Exploration}. Springer, 2016.
\bibitem[LF26a]{LF26a} Bruno Le Floch. \emph{Implication semilattice of 990 quasigroup equational laws}. arXiv:2603.29909, 2026. \url{https://arxiv.org/abs/2603.29909}
\bibitem[LF26b]{LF26b} Bruno Le Floch. \emph{Lattice of 456 semigroup varieties from equations of order up to 4}. arXiv:2609.39872, 2026. \url{https://arxiv.org/abs/2609.39872}
\end{thebibliography}

\end{document}
